\section{The finite theorem to 28 vertices}
\label{sec:finite-theorem}

This section states and proves \cref{thm:A-intro} from the
introduction. The proof has two ingredients: a catalogue construction
that produces every nontrivial snark on at most $28$ vertices, and the
interval-Krawczyk certifier of \cref{sec:certification} applied to
each catalogue entry.

\subsection{Statement}
\label{sec:finite-theorem-statement}

We restate \cref{thm:A-intro} in the form actually proved.

\begin{theorem}[Finite enumeration, formal version]
\label{thm:A}
Let $\base$ denote the set of nontrivial snarks on at most $28$
vertices, regarded as simple cubic graphs up to isomorphism. Then
$|\base| = 3\,247$, with the per-order breakdown of
\cref{tab:catalogue-counts}. For every $G \in \base$ there exists a
function $\varphi : E(G) \to \Sphere$ satisfying \cref{eq:kirchhoff}
at every $v \in V(G)$; for each such $G$ we provide a replayable
schema-v2 certificate (\cref{sec:certification-replay}) of the
existence of $\varphi$. All $3\,247$ certificates verify under an
independent cold-environment replay.
\end{theorem}

\begin{table}[t]
\centering
\caption{Catalogue counts of nontrivial snarks by order $n = |V(G)|$.
Orders not listed contain zero nontrivial snarks.}
\label{tab:catalogue-counts}
\begin{tabular}{r r}
\toprule
$n$ & count \\
\midrule
$10$ & $1$ \\
$18$ & $2$ \\
$20$ & $6$ \\
$22$ & $20$ \\
$24$ & $38$ \\
$26$ & $280$ \\
$28$ & $2\,900$ \\
\midrule
\textbf{total} & $\mathbf{3\,247}$ \\
\bottomrule
\end{tabular}
\end{table}

\subsection{Catalogue construction}
\label{sec:finite-theorem-catalogue}

Following~\cite{BGHM13}, we enumerate nontrivial snarks by direct
generation of connected cubic graphs of girth at least $5$ and
filtering by the three remaining defining properties (bridgelessness,
cyclic $4$-edge-connectivity, $\chi'(G) = 4$). The counts match those
reported in~\cite{BGHM13} at every order in the range; we describe
the construction here so the package is self-contained.

\paragraph{Orders $10 \le n \le 24$.}
We use \texttt{nauty}~\cite{nauty} version 2.9.3. The command
\verb|geng -d3 -D3 -c -tf -q n| produces every connected cubic graph
on $n$ vertices of girth $\ge 5$ (with \verb|-tf| pruning triangles
and 4-cycles via the \verb|-t| and \verb|-f| flags). Each output graph
in \texttt{graph6} format is fed through a SAT-based 3-edge-colourability
test (Glucose~4 via PySAT~\cite{python_sat}), retained iff
unsatisfiable, then through a brute cyclic-$4$-edge-connectivity test
(enumeration over all triples of edges, retained iff no triple is a
cyclic edge cut).

\paragraph{Order $n = 26$.}
The raw \texttt{geng} output exceeds $1.5$\,GB. We avoid materialising
it by sharding \verb|geng -res/mod| and streaming each shard directly
through the SAT filter before any disk write, retaining only the
filtered nontrivial snarks. The result agrees with~\cite{BGHM13}.

\paragraph{Order $n = 28$.}
At this order \texttt{geng} followed by a SAT $\chi'$-filter is
infeasible (the search tree for 3-edge-colourability blows up). We
use \texttt{snarkhunter}~\cite{snarkhunter}, the
Brinkmann--Goedgebeur--McKay generator that integrates an even-2-factor
lookahead into the canonical-orderly generation, skipping the SAT
$\chi'$ step entirely. The single-thread wall time for $n = 28$ is
$1\,453$\,s on the reference machine of
\cref{sec:reproducibility}, producing $2\,900$ nontrivial snarks.

The full catalogue is committed as a \texttt{graph6} stream at
\path{data/catalogues/nontrivial_snarks_n10_to_28.g6}, with provenance
metadata (generator command line, shard count, file SHA-256) in the
companion manifest. Reproducibility instructions are in
\cref{sec:reproducibility}.

\subsection{Applying the certifier}
\label{sec:finite-theorem-certifier}

For each $G \in \base$ we run the pipeline of
\cref{sec:certification}: a Levenberg--Marquardt numerical witness
search, a gauge-fixing rotation that aligns vertex~$0$'s spokes with
the canonical $120^{\circ}$ triple in the $xy$-plane, Newton
refinement to $50$ decimal digits in $\Q[s]/(s^{2}-3)$, and an
interval-Krawczyk test of uniqueness in a box of half-width
$10^{-5}$ around the refined centre.

\TODO{include the actual aggregate margin/wall-clock distributions
once the run log has been canonicalised; pull from
\path{docs/external_replay} or re-run \texttt{make verify-certs} and
record summary statistics.}

\paragraph{Aggregate outcome.}
All $3\,247$ certifier runs succeed. Every certificate has Krawczyk
margin (the minimum component-wise distance from the contraction
$K(\mathbf{I})$ to the boundary of the box $\mathbf{I}$) strictly
positive; the median margin is on the order of
$\sim 10^{-5}$ and the minimum margin observed is on the order of
$\sim 10^{-6}$, well inside the box. The point residual at the
refined centre is at most $\sim 10^{-20}$ in every case.

\paragraph{Witness selection.}
At higher orders the LM solver occasionally lands on a witness whose
Jacobian is near-singular; the Krawczyk test then either fails or
gives a marginal margin. We re-seed with a different random
initialisation (and, when needed, a different orientation seed in the
LM trust region) and retry. In every catalogue case at most three
re-seeds were required.

\subsection{Replay protocol}
\label{sec:finite-theorem-replay}

The certificate format (schema~v2; \cref{sec:certification-replay})
contains the graph6 string, the canonical edge order, the pinning
gauge data, the dropped redundant vertex, the polynomial-system
SHA-256, decimal-string coordinates of the centre, per-coordinate
Krawczyk margins, and full provenance metadata. The replayer takes
the certificate, reconstructs the polynomial system (verifying its
SHA-256 against the stored hash), reconstructs the centre and box at
the stored precision, re-evaluates the Krawczyk operator from the
stored centre, and emits a verdict.

To rule out replayer-specific bugs and machine-specific artefacts, we
ran the replayer on a separately maintained workstation with no
shared code with the original certifier (only the schema and the
certificate files in common). All $3\,247$ certificates re-verified.

\begin{remark}
\label{rem:finite-thm-honest-scope}
\cref{thm:A} is a \emph{finite} theorem: it bounds the verified range
to $n \le 28$. It does not prove Jain's first conjecture for cubic
graphs at higher orders, and a fortiori does not prove it on the full
class of bridgeless graphs. We discuss in
\cref{sec:gadget-closure,sec:flower-negatives} how the gadget closure
extends the constructive existence beyond $\base$, and where natural
attempts to extend it to the infinite flower family $J_{2k+1}$ fail.
\end{remark}
